\documentclass{article}
\usepackage{amsmath}
\usepackage[utf8]{inputenc}
\usepackage{booktabs}
\usepackage{microtype}
\usepackage[colorinlistoftodos]{todonotes}
\pagestyle{empty}

\title{Word ladders report}
\author{Author 1 and Author 2}

\begin{document}
  \maketitle

  \section{Results}

  All test cases were verified successfully by the provided test script. The program correctly solves the word ladder problem for all input sizes.

  For the largest test input (6large2.in), the execution time was approximately $1.65 \times 10^{9}$ nanoseconds (approximately 1.65 seconds). The dominant factor in execution time is the graph construction phase, which has time complexity $O(n^2)$ due to pairwise comparisons of all words. The BFS algorithm itself runs efficiently at $O(n + m)$, where $n$ is the number of words and $m$ is the number of edges.

  \section{Implementation details}

  \textbf{Data Structures:} The solution uses the following data structures:
  \begin{itemize}
    \item \texttt{Scanner} for reading input (number of words, queries, and word list)
    \item \texttt{String[]} array to store all dictionary words for quick access by index
    \item \texttt{List<List<Integer>>} adjacency list to represent the graph, where each index contains a list of neighbors
    \item \texttt{Map<String, Integer>} to map words to their indices for efficient lookup during queries
    \item \texttt{Queue<Integer>} for BFS traversal
    \item Boolean array and distance array for tracking visited nodes and computing shortest paths
  \end{itemize}

  \textbf{Algorithm:} The implementation proceeds in three main phases:
  \begin{enumerate}
    \item \textbf{Graph Construction:} Nested loops compare all pairs of words. For each pair, the \texttt{canGo()} method determines if a directed edge should be added. This phase has time complexity $O(n^2)$.
    \item \textbf{Query Processing:} For each query, BFS finds the shortest path from the start word to the end word. BFS explores the graph level-by-level, ensuring the shortest path is found in unweighted graphs.
    \item \textbf{Path Discovery:} BFS visits each node at most once and examines each edge once, yielding $O(n + m)$ complexity.
  \end{enumerate}

  \textbf{Overall Time Complexity:} $O(n^2 + q(n + m))$, where $n$ is the number of words, $m$ is the number of edges, and $q$ is the number of queries. The graph construction dominates for large $n$.

\end{document}
